% Options for packages loaded elsewhere
\PassOptionsToPackage{unicode}{hyperref}
\PassOptionsToPackage{hyphens}{url}
\documentclass[
]{article}
\usepackage{xcolor}
\usepackage{amsmath,amssymb}
\setcounter{secnumdepth}{-\maxdimen} % remove section numbering
\usepackage{iftex}
\ifPDFTeX
  \usepackage[T1]{fontenc}
  \usepackage[utf8]{inputenc}
  \usepackage{textcomp} % provide euro and other symbols
\else % if luatex or xetex
  \usepackage{unicode-math} % this also loads fontspec
  \defaultfontfeatures{Scale=MatchLowercase}
  \defaultfontfeatures[\rmfamily]{Ligatures=TeX,Scale=1}
\fi
\usepackage{lmodern}
\ifPDFTeX\else
  % xetex/luatex font selection
\fi
% Use upquote if available, for straight quotes in verbatim environments
\IfFileExists{upquote.sty}{\usepackage{upquote}}{}
\IfFileExists{microtype.sty}{% use microtype if available
  \usepackage[]{microtype}
  \UseMicrotypeSet[protrusion]{basicmath} % disable protrusion for tt fonts
}{}
\makeatletter
\@ifundefined{KOMAClassName}{% if non-KOMA class
  \IfFileExists{parskip.sty}{%
    \usepackage{parskip}
  }{% else
    \setlength{\parindent}{0pt}
    \setlength{\parskip}{6pt plus 2pt minus 1pt}}
}{% if KOMA class
  \KOMAoptions{parskip=half}}
\makeatother
\usepackage{longtable,booktabs,array}
\newcounter{none} % for unnumbered tables
\usepackage{calc} % for calculating minipage widths
% Correct order of tables after \paragraph or \subparagraph
\usepackage{etoolbox}
\makeatletter
\patchcmd\longtable{\par}{\if@noskipsec\mbox{}\fi\par}{}{}
\makeatother
% Allow footnotes in longtable head/foot
\IfFileExists{footnotehyper.sty}{\usepackage{footnotehyper}}{\usepackage{footnote}}
\makesavenoteenv{longtable}
\setlength{\emergencystretch}{3em} % prevent overfull lines
\providecommand{\tightlist}{%
  \setlength{\itemsep}{0pt}\setlength{\parskip}{0pt}}
\usepackage{hyperxmp}
\usepackage{bookmark}
\IfFileExists{xurl.sty}{\usepackage{xurl}}{} % add URL line breaks if available
\urlstyle{same}
\hypersetup{
  pdftitle={Most infrastructure layers are symptoms of the persistence model: a construct for auditing production stacks},
  pdfauthor={Alvaro Rivera},
  pdfkeywords={\xmpquote{infrastructural symptom}, \xmpquote{design
theory}, \xmpquote{persistence model}, \xmpquote{compensatory
layer}, \xmpquote{journaled systems}, \xmpquote{actor-native
architecture}, \xmpquote{Redis}, \xmpquote{DBMS-centric
architecture}, \xmpquote{embedded deployment}, \xmpquote{puppeteer
framework}, \xmpquote{accidental category}, \xmpquote{server-role}},
  hidelinks,
  pdfcreator={LaTeX via pandoc}}

\title{Most infrastructure layers are symptoms of the persistence model:
a construct for auditing production stacks}
\author{Alvaro Rivera}
\date{2026-05-15}

\begin{document}
\maketitle
\begin{abstract}
Production software systems routinely include layers such as caches,
object-relational mappers, message queues, distributed locks,
application server frameworks, and orchestration clusters. These layers
are treated as standard components of non-trivial deployments, yet they
are rarely examined under a common explanation for why they exist.

This paper argues that many of these layers do not arise from structural
requirements of the problem being solved, but from structural
deficiencies of the underlying persistence model. It names this property
\emph{infrastructural symptom}. A layer exhibits infrastructural symptom
when it compensates for a defect of the persistence model and loses
justification when that defect is removed. Two conditions ---
compensation and dissolution --- formalize this distinction, and a
forensic procedure operationalizes it for any (layer, use-case) pair
without requiring adoption of a specific alternative.

The paper presents a journaled actor-native system as an existence proof
in which the canonical symptoms dissolve. Four mechanisms --- locality,
journal density, journal-as-causal-substrate, and self-contained
substrate --- account for the disappearance of seven common categories
of compensating infrastructure. The construct extends through a ladder
of additional layers to a limit case on minimal hardware, as an
architectural permission rather than an operational demonstration.

The contribution is a design theory contribution in the sense of Hevner,
March, Park, and Ram (2004): a recognition construct that makes visible
which parts of a stack compensate for the model rather than serve the
problem.
\end{abstract}

\section{Most infrastructure layers are symptoms of the persistence
model: a construct for auditing production
stacks}\label{most-infrastructure-layers-are-symptoms-of-the-persistence-model-a-construct-for-auditing-production-stacks}

\subsection{TL;DR}\label{tldr}

Many layers taken for granted in production systems --- caches, ORMs,
message queues, distributed locks, application servers, orchestration
clusters --- do not exist because the problem requires them, but because
the underlying persistence model does.

This paper names that property \emph{infrastructural symptom}, defines
the two conditions that distinguish it from a structural requirement,
provides a forensic procedure to audit any layer under those conditions,
and shows a journaled actor-native system in which the canonical
symptoms disappear.

The construct is diagnostic, not prescriptive: it makes visible which
parts of a stack compensate for the model rather than serve the problem.

\subsection{Claims this paper makes}\label{claims-this-paper-makes}

\begin{enumerate}
\def\labelenumi{\arabic{enumi}.}
\tightlist
\item
  Some infrastructural layers in production deployments exist to
  compensate for structural deficiencies of the underlying persistence
  model rather than for structural requirements of the problem being
  solved (§3).
\item
  The property naming this distinction --- \emph{infrastructural
  symptom} --- is operationalized by two conditions: compensation (the
  layer compensates for an identifiable model defect) and dissolution
  (the layer loses justification when the defect is removed) (§3).
\item
  A forensic procedure operationalizes the two conditions and produces a
  binary verdict per (layer, use-case) pair (§5).
\item
  A journaled actor-native model instantiates a system in which four
  mechanisms --- locality, journal density, journal-as-causal-substrate,
  self-contained substrate --- dissolve the seven canonical categories
  of compensating infrastructure (§4.1).
\item
  The Redis case admits the procedure with per-use-case discrimination:
  six use-cases dissolve as symptoms; two remain structural (§4.2).
\item
  The ladder of compensated layers is open-ended; the procedure extends
  to backup tooling, service mesh, distributed-tracing/APM, and future
  emergent layers (§6).
\item
  The model permits embedded-hardware deployment as a limit case --- an
  architectural permission, not an operational proof (§7).
\end{enumerate}

\subsection{1. Introduction}\label{introduction}

This paper makes a design theory contribution in the sense of Hevner,
March, Park, and Ram (2004). It identifies and names a structural
pattern that prior literature has documented case by case without
articulating as one (the construct); derives the conditions under which
the pattern dissolves (the principles); and presents an instantiation in
which those conditions hold and the pattern is absent (the existence
proof). The genre is design science research: evidence is presented in
the form of a working artifact rather than a controlled experiment. The
contribution is conceptual; the instantiation confirms that the pattern
is contingent rather than structurally necessary, not that any specific
system is the only way to dissolve it.

Production software systems are routinely deployed alongside layers of
supporting infrastructure: in-memory caches, object-relational mappers,
message queues, distributed locks, application server frameworks,
orchestration clusters, full operating systems. These layers are treated
as standard components of any non-trivial deployment. Architectural
discussions assume their presence by default; tutorials and reference
architectures position them as foundational; engineering teams budget
for them, staff them, and monitor them. Their absence is treated as a
property of toy systems.

This paper names a property that some --- not all --- of these layers
share. We call this property \emph{infrastructural symptom}: a layer is
an infrastructural symptom when it exists to compensate for a structural
deficiency in the underlying persistence model rather than for a
structural requirement of the problem being solved. The distinction is
routinely elided; the two cases are treated alike in architecture, in
tooling, and in engineering culture. Once the property is named, each
layer in a given stack becomes auditable under the question: is this
compensating for something, and if so, for what?

Section 2 situates the paper among prior work. Section 3 defines the
construct formally and derives the two conditions that distinguish
infrastructural symptom from structural requirement. Section 4 presents
the instantiation: a journaled actor-native system in which the
canonical layers --- caches, glue queues, locks, application server
scaffolding --- are absent not by omission but by structural
irrelevance. Section 5 specifies a forensic procedure for auditing a
stack under the construct. Section 6 generalizes from the canonical case
to a chain of compensated layers, from in-memory cache up through full
operating system. Section 7 examines the limit case, where the chain
collapses to a domain artifact running on minimal hardware. Section 8
examines where the construct does not apply.

\subsection{2. Related work}\label{related-work}

The genre of this paper is design science research as articulated by
Hevner, March, Park, and Ram (2004). A design theory paper identifies a
construct, derives the principles under which the construct's predicate
operates, and presents an instantiation that demonstrates the construct
is realizable. The contribution is conceptual; the instantiation is the
existence proof, not the substance of the claim.

Critique of DBMS-centric architecture as the dominant paradigm for
production systems has accumulated in the literature without naming the
construct presented here. Stonebraker and Çetintemel (2005) observe that
the relational database has become a one-size-fits-all hammer applied to
use-cases for which it is structurally ill-suited, and predict the
proliferation of specialized stores. Kleppmann (2017) catalogs the
patterns by which contemporary deployments combine caches, queues, ORMs,
and search indices around a relational core, documenting their pairwise
interactions and operational costs. Helland's \emph{Life beyond
Distributed Transactions} (2007) makes a parallel observation about the
proliferation of coordination machinery under distributed-transaction
semantics. His \emph{Immutability Changes Everything} (2016) extends a
related argument to the substrate itself: append-only representation
dissolves much of that machinery by construction. The present paper
names what these critiques document piecewise --- the unifying property
of compensation --- and operationalizes the discrimination.

The instantiation presented in §4 is constructed from three lines of
prior work. The actor model originates with Hewitt (1973), articulating
computation as message passing between isolated entities. Event sourcing
as a persistence pattern is associated with Fowler (2005), which
characterizes the program's history as the durable artifact rather than
a snapshot of current state. Command Query Responsibility Segregation,
articulated by Young (2010), separates the read and write paths so they
no longer compete for the same model. Vernon (2015) brings these three
together as practical patterns in \emph{Reactive Messaging Patterns with
the Actor Model}; that synthesis is the immediate ancestor of the
present instantiation. The journaled actor-native system used as the
existence proof here combines these threads; it does not invent them.

This paper is the sixth in a series. The series develops the construct's
foundations across five preceding papers. Paper 1 introduces porosity as
the representational defect that this paper's construct identifies in
its compensation form. Paper 2 develops program-value separability and
the externalized-parameters precondition that supports journal density
as a §4.1 mechanism. Paper 3 introduces Reactions and the pragmatic
partition between work-due-before-responding and deferred work, which
§4.1 invokes as the journal-as-causal-substrate mechanism. Paper 4
introduces the Tell primitive and cross-actor causal continuity, used in
§4.1 for the same mechanism. Paper 5 develops the journal as substrate
for deployment, replication, backup, and offline operation, which §6
references as the substrate-mediated lifecycle that dissolves
application-server scaffolding. Paper 7 carries the same discrimination
upward from infrastructural layers to architectural roles; §9 introduces
this extension.

The contribution unique to this paper is conceptual. The construct
\emph{infrastructural symptom} names a property that prior work
documents in pieces but does not unify. The two conditions ---
compensation and dissolution --- formalize the property. The forensic
procedure of §5 operationalizes it without requiring adoption of any
particular alternative. The ladder of §6 and the limit case of §7 extend
the construct's reach. The instantiation of §4 demonstrates that the
construct's predicate dissolves under the conditions; the demonstration
is not the contribution.

\subsection{3. The construct: infrastructural
symptom}\label{the-construct-infrastructural-symptom}

We state the construct formally before deriving its two conditions and
the apparatus that operationalizes them.

\textbf{Formally.} An infrastructural symptom is a layer of
infrastructure that exists to compensate for a structural deficiency in
the underlying persistence model rather than for a structural
requirement of the problem being solved.

The construct operates at the level of purpose, not at the level of
product. A given physical layer in a deployment may serve multiple
purposes; the construct discriminates among those purposes one at a
time. The unit of analysis is the use-case, not the layer.

A layer is an infrastructural symptom for a given use-case when both of
two conditions hold.

\textbf{The compensation condition.} The layer's purpose for that
use-case is traceable to a specific deficiency of the underlying
persistence model. The deficiency must be identifiable: not ``we use
this layer because it is fast,'' but ``we use this layer because the
persistence model exhibits property P, and this layer compensates for
P.'' When the compensation condition holds, asking what would happen if
P were not present? yields a non-trivial answer.

\textbf{The dissolution condition.} When that deficiency is removed ---
by adopting a persistence model in which P does not arise --- the layer
loses justification for that use-case. The deficiency's absence is
sufficient to make the layer redundant for that purpose; nothing else
about the problem domain, workload, or operational context needs to
change.

Both conditions must hold. Compensation without dissolution describes a
genuine patch whose value survives model evolution. Dissolution without
compensation is not constructible: a layer cannot dissolve under a model
change unless its purpose was tied to the model in the first place.

The construct is not a universal accusation against infrastructure. Many
layers in production systems exist as structural requirements --- they
serve purposes rooted in the problem domain, independent of any
persistence model. Three categories are typical:

\begin{itemize}
\tightlist
\item
  \textbf{External-system integration.} Layers mediating communication
  with systems outside the boundary --- payment gateways, partner APIs,
  sensor networks --- compensate for no model deficiency. The
  integration problem remains under any persistence model.
\item
  \textbf{Cryptographic and protocol enforcement.} TLS termination,
  request signing, OAuth verification address security or protocol
  requirements rooted in the problem domain. No persistence-model change
  makes encryption optional.
\item
  \textbf{Genuinely expensive computation.} Analytics engines,
  machine-learning inference services, video encoders exist because the
  computation itself is intrinsically expensive. That cost does not
  change under a different persistence model.
\end{itemize}

The boundary between symptom and structural requirement is not a
property of the layer as a product but of the purpose for which it is
used. The same caching system, deployed to absorb read-vs-write
contention at the persistence layer, instantiates an infrastructural
symptom. Deployed instead to memoize a deliberately expensive
computation, it instantiates a structural requirement. The discriminator
operates per use-case.

Seven categories of layer are encountered regularly in production
deployments and admit discrimination under the two conditions. The third
column describes, in model-property terms rather than system terms, what
a persistence model satisfying the two conditions would offer instead.

\textbf{Table 1.} Discriminator: seven canonical categories of
compensating infrastructure under the two conditions of §3.

{\def\LTcaptype{none} % do not increment counter
\begin{longtable}[]{@{}
  >{\raggedright\arraybackslash}p{(\linewidth - 4\tabcolsep) * \real{0.3333}}
  >{\raggedright\arraybackslash}p{(\linewidth - 4\tabcolsep) * \real{0.3333}}
  >{\raggedright\arraybackslash}p{(\linewidth - 4\tabcolsep) * \real{0.3333}}@{}}
\toprule\noalign{}
\begin{minipage}[b]{\linewidth}\raggedright
Layer
\end{minipage} & \begin{minipage}[b]{\linewidth}\raggedright
Underlying defect compensated
\end{minipage} & \begin{minipage}[b]{\linewidth}\raggedright
What a model satisfying the conditions offers instead
\end{minipage} \\
\midrule\noalign{}
\endhead
\bottomrule\noalign{}
\endlastfoot
In-memory cache & Reads contend with writes at the persistence layer;
reconstructing state from storage is expensive & State local to the unit
of computation; no contention between read and write paths \\
Object-relational mapper & Domain objects do not align with relational
rows; mapping is opaque and lossy & Domain objects are the unit of
persistence; no translation layer \\
Message queue used as glue & Components cannot directly invoke deferred
work without losing causal record & Deferred work is journal-recorded
and replayable; queues are not the substrate of causality \\
Distributed lock & Multiple writers contend for the same logical entity
across processes & A single point of authority per logical entity;
serialization is intrinsic, not externalized \\
Application server framework & The persistence model leaks into
application code; lifecycle plumbing must be added & The domain is the
application; lifecycle is a property of the substrate, not added
scaffolding \\
Orchestration cluster as coordination substrate & Stateful coordination
across instances requires external choreography & State has location;
coordination is between actors, not between containers \\
Full operating system & The runtime, the framework, and the application
require disparate substrates & A minimal substrate suffices to run the
domain artifact \\
\end{longtable}
}

The seven rows are not exhaustive. Backup tooling, service mesh, and
distributed-tracing frameworks admit similar analysis and are discussed
in Section 6. Whether such a persistence model exists, and at what cost,
is the subject of Section 4.

The construct is a recognition construct, not a prescription. It enables
auditing existing stacks under the two conditions; it does not specify
how to build a model in which infrastructural symptoms do not arise.
Construction is the work of Section 4, which exhibits one such model,
and of Section 5, which specifies the forensic procedure that produced
the discriminator table.

The value of the construct is diagnostic. A practitioner who does not
adopt the instantiation of Section 4 can nevertheless use the construct
to determine which layers in their stack are tied to their persistence
model and which are tied to their problem domain. That visibility is
independent of the choice to act on it.

\subsection{4. Instantiation}\label{instantiation}

\subsubsection{4.0 Genealogy}\label{genealogy}

Puppeteer's earliest layer dates to 2005, when a persistence library
internally called \emph{autopersistencia} adopted a single design rule:
the classes representing the domain would carry no persistence concerns.
Persistence would be discovered by reflection over the class hierarchy;
storage details lived elsewhere. The rule was modest in scope and
pragmatic in motivation. Its consequence was structural: the domain
became a clean substrate, unobstructed by the technical aspects that
conventional architectures embed inside it.

At no point in this evolution was there an intention to eliminate
caches, queues, locks, or application scaffolding; those layers were
assumed to belong in any serious system. Their later absence was not a
design objective but an empirical outcome. This distinction matters for
the argument of this paper: the construct is not dissolved here by
architectural ideology but as a byproduct of unrelated design
constraints that happened to satisfy the conditions derived in §3.

The natural follow-up question was: if domain classes carry no
persistence concerns, where does state live? The answer that emerged
across the following years became the framework's first generation.
State lives in a journal, and the journal records the program itself ---
not serialized events, not row deltas, but the DSL script that produced
the state, replayable in order. The narrative of execution, not the
photograph of memory, became the unit of persistence --- a property
whose absence, in the terms of §3, is the deficiency that caches, ORMs,
and queue-glue exist to compensate for in DBMS-centric architectures.

The framework's second generation added four properties, each developed
in preceding papers of this series. Parameters were externalized,
allowing programs to be compiled, cached, and recorded as references
rather than literal scripts (Paper 2). Reactions were introduced as the
developer-controlled partition between work due before responding and
work deferred to placement (Paper 3). The Tell primitive established
cross-actor causality recorded in the journal rather than orchestrated
externally (Paper 4). The journal itself became the catalog of the
actor's declarative vocabulary, dissolving the lateral storage of action
definitions that the framework had carried since V1.

The framework was not designed to dissolve the construct named in this
paper. The construct was identified retrospectively, after the chain of
architectural decisions described above had converged on a system in
which most of the canonical layers had become unnecessary. The work of
this paper is, in that sense, a forensic reading of an existing
artifact, not an engineering plan derived from a principle. The
principle was a property of the artifact before it was a property of the
literature.

\subsubsection{4.1 The instantiation under the two
conditions}\label{the-instantiation-under-the-two-conditions}

Four mechanisms in the journaled actor-native model account for the
dissolution of the seven canonical categories of §3. A single mechanism
may dissolve more than one category, and the in-memory cache dissolves
under two mechanisms because §3 records two distinct defects compensated
by that layer. Each mechanism decomposes into several sub-mechanisms ---
properties of the model that, taken together, exhibit the umbrella
behavior. Table 2 correlates each mechanism with its sub-mechanisms, the
§3 defects it dissolves, and the canonical layers that become symptoms
under it.

\textbf{Table 2.} Mechanisms in the journaled actor-native model mapped
to the §3 defects they dissolve and the canonical layers that become
symptoms.

{\def\LTcaptype{none} % do not increment counter
\begin{longtable}[]{@{}
  >{\raggedright\arraybackslash}p{(\linewidth - 4\tabcolsep) * \real{0.3333}}
  >{\raggedright\arraybackslash}p{(\linewidth - 4\tabcolsep) * \real{0.3333}}
  >{\raggedright\arraybackslash}p{(\linewidth - 4\tabcolsep) * \real{0.3333}}@{}}
\toprule\noalign{}
\begin{minipage}[b]{\linewidth}\raggedright
Mechanism
\end{minipage} & \begin{minipage}[b]{\linewidth}\raggedright
§3 defect compensated
\end{minipage} & \begin{minipage}[b]{\linewidth}\raggedright
Canonical layer that becomes symptom
\end{minipage} \\
\midrule\noalign{}
\endhead
\bottomrule\noalign{}
\endlastfoot
\textbf{Locality}~~• privacy of state~~• actor-as-serializer~~•
placement & Reads contend with writes at the persistence layerMultiple
writers contend for the same logical entity & In-memory cache
(contention defect)Distributed lock \\
\textbf{Journal density}~~• journal-as-program (homoiconic)~~•
replay-as-reconstruction~~• domain-class-as-unit-of-persistence &
Reconstructing state from storage is expensiveDomain objects do not
align with relational rows & In-memory cache (reconstruction
defect)Object-relational mapper \\
\textbf{Journal as causal substrate}~~• Reactions recorded in journal~~•
Tell recorded in journal~~• causality boundary is the journal &
Components cannot directly invoke deferred work without losing causal
recordStateful coordination across instances requires external
choreography & Message queue used as glueOrchestration cluster \\
\textbf{Self-contained substrate}~~• domain artifact is the
application~~• substrate carries lifecycle, persistence, dispatch~~•
journal-mediated release handoff~~• minimal OS surface & The persistence
model leaks into application codeThe runtime, framework, and application
require disparate substrates & Application server frameworkFull
operating system \\
\end{longtable}
}

\textbf{Locality.} State is private to the actor that owns it; the actor
processes commands and queries serially against its own memory. There is
no shared store from which concurrent reads must be isolated from
concurrent writes. Two of the §3 defects vanish mechanically under this
property: the read-versus-write contention that motivates in-memory
caches, and the inter-process contention that motivates distributed
locks. The compensation condition is vacuous for both. Placement is the
developer-controlled partition treated in Paper 3.

\textbf{Journal density.} The journal does not record serialized
snapshots or row-by-row deltas; it records the program --- the DSL
script --- that produced the state. Reconstruction is replay, not
deserialization; replay of a compact script is cheap. The domain class
is the unit of persistence, with no intermediate row to map. Both §3
defects compensated by caches and ORMs --- reconstruction expense and
object-row impedance --- are dissolved mechanically. The mechanism is
treated as compilation in Paper 2 and as homoiconic representation in
Paper 1. Implementation: \texttt{Puppeteer/EventSourcing/DB/Diary.cs}.

\textbf{Journal as causal substrate.} Pragmatic deferral, developed in
Paper 3 as Reactions, records work the verb does not perform before
responding as a journal entry, with the causal link to the originating
event preserved. Cross-actor continuity, developed in Paper 4 as Tell,
records dispatch from one actor to another in the journal as well. The
journal becomes the boundary of causality. In DBMS-centric systems,
causality crosses the persistence boundary and must be reconstructed
with message queues used as glue and orchestration clusters used as
coordination substrates; here, causality never leaves the journal.
Implementation: the \texttt{reactions} field in
\texttt{Puppeteer/EventSourcing/ActorHandler.cs:38}.

\textbf{Self-contained substrate.} The runtime is a small loader; the
application is the domain artifact. Lifecycle, persistence, and dispatch
are properties of the substrate, so no external process is required to
host or coordinate the application artifact. Release handoff and rolling
deployment, handled in DBMS-centric deployments by application server
orchestration, are journal-mediated here and are the subject of Paper 5.
The application server framework as a category has no equivalent role;
the operating system surface is reduced to the minimum required, a limit
case treated in §7.

Together, these four mechanisms account for the dissolution of the seven
canonical categories of §3 mechanically, not by intent. None was added
to the model with the construct of this paper in mind; each was
developed for the reasons enumerated in §4.0.

\subsubsection{4.2 The Redis case in
detail}\label{the-redis-case-in-detail}

Redis occupies a distinctive position in the contemporary infrastructure
landscape: a single product whose presence is taken for granted across
deployments that are otherwise architecturally dissimilar. The diversity
of its typical use-cases --- cache, session store, pub/sub, queue,
distributed lock, leaderboard, rate limiter --- does not reflect the
diversity of a single problem domain. It reflects the diversity of the
pressures that DBMS-centric architectures present to application
designers, each of which Redis conveniently absorbs in a familiar form.
This makes Redis the canonical case for inspection under the construct
of §3. The analysis that follows is not, however, about Redis as a
product; it is about the defects of the underlying persistence model
that Redis exists to absorb --- defects that other compensating layers
also address in variant forms, treated briefly after the canonical
table.

Table 3 applies the two conditions of §3 to the most common Redis
use-cases, tracing each to the specific defect it compensates and the
mechanism in §4.1 that dissolves it.

\textbf{Table 3.} Application of the construct to canonical Redis
use-cases.

{\def\LTcaptype{none} % do not increment counter
\begin{longtable}[]{@{}
  >{\raggedright\arraybackslash}p{(\linewidth - 4\tabcolsep) * \real{0.3333}}
  >{\raggedright\arraybackslash}p{(\linewidth - 4\tabcolsep) * \real{0.3333}}
  >{\raggedright\arraybackslash}p{(\linewidth - 4\tabcolsep) * \real{0.3333}}@{}}
\toprule\noalign{}
\begin{minipage}[b]{\linewidth}\raggedright
Redis use-case
\end{minipage} & \begin{minipage}[b]{\linewidth}\raggedright
§3 defect compensated
\end{minipage} & \begin{minipage}[b]{\linewidth}\raggedright
Native response in journaled actor-native
\end{minipage} \\
\midrule\noalign{}
\endhead
\bottomrule\noalign{}
\endlastfoot
In-process cache for database results & Reads contend with writes at the
persistence layer; reconstructing state from storage is expensive &
Actor state is the cache (§4.1 Locality + Journal density) \\
Session storage & Stateless application servers cannot retain user state
across requests & The user actor retains session naturally as part of
its state (§4.1 Locality + Self-contained substrate) \\
Pub/sub among internal components & No causal record of cross-component
events; coordination must be reconstructed & Reactions record deferred
work in the journal with causal links preserved (§4.1 Journal as causal
substrate) \\
Message queue for deferred work & Components cannot directly invoke
deferred work without losing causal record & Reactions and Tell record
causation in the journal (§4.1 Journal as causal substrate) \\
Distributed lock for shared entity & Multiple processes write the same
logical entity across boundaries & The actor owns its entity;
serialization is intrinsic, not externalized (§4.1 Locality) \\
Leaderboard / sorted set & Real-time aggregation against the persistence
layer is expensive & Query against actor state or a follower; replay is
cheap (§4.1 Locality + Journal density) \\
\end{longtable}
}

Each row satisfies both the compensation and dissolution conditions of
§3. The pub/sub row covers internal-component coordination; the
external-systems variant is structural and is treated below.

Two Redis use-cases do not satisfy the dissolution condition and are
therefore structural requirements under the construct. The first is
pub/sub directed at external systems: webhooks to third-party consumers,
event distribution to partners, broadcast to subscribers outside the
operator's boundary. The external system does not dissolve under any
persistence-model change; the distribution problem persists. The second
is rate limiting applied to external API consumption: per-customer
quotas, third-party API throttling, abuse mitigation against external
traffic. The rate limit is a property of the external relationship, not
of the internal model. Redis, or any comparable substrate, remains the
appropriate tool for both.

Other compensating layers in the DBMS-centric ecosystem exhibit variants
of the same pattern. Memcached shares the cache analysis: locality and
journal density make in-process caches mechanically redundant.
Distributed coordinators such as ZooKeeper share the lock and
coordination analysis: the actor model makes the underlying contention
vacuous. Message brokers such as RabbitMQ share the queue analysis:
Reactions and Tell record causation in the journal.

Event-storage products such as EventStore DB and Kafka deserve a
separate observation. Both adopt a journaled posture --- they record
events rather than serialized snapshots, and they offer replay. But they
instantiate the journal as an external infrastructure layer that the
application consumes through a client API, not as the substrate of the
application itself. Causality is recorded but remains external to the
unit of computation; coordination between event production, event
consumption, and application state still requires application-side glue.
The product captures the right idea and externalizes it, leaving the
application code in the compensating position. Under the construct,
these products do not dissolve the symptom; they relocate it.

A second observation about these compensating layers is that they do not
absorb their compensation freely. Each requires the application to carry
the boilerplate of its use: cache invalidation logic, lock acquisition
and release with retry and fencing semantics, serialization between
domain objects and external representations, key naming and TTL
discipline, consumer offset management and rebalancing handlers, schema
registry coordination, idempotency keys for at-least-once delivery. The
application code carries the porosity imposed by RDBMS representations
(Paper 1) as its baseline; the compensating layer adds a second
porosity, this time imposed by the layer itself: its own API surface,
lifecycle, and consistency contracts. Domain logic interleaves with
both. Under the instantiation, the compensating layer and the
boilerplate that the application carries to use it are both absent. The
actor's code addresses the problem domain; the substrate addresses the
rest.

Redis as a product is not what the construct identifies; nor is any
single compensating layer treated above. What the construct identifies
is the architectural pattern in which a persistence model exposes
defects and successive layers --- caches, brokers, coordinators, even
externalized event stores --- accumulate to absorb them. Replacing one
such layer while retaining the model that produces the others does not
satisfy the construct's conditions. The unit of change is the
persistence model itself, not any individual layer.

\subsection{5. Forensic procedure}\label{forensic-procedure}

The construct of §3 enables discrimination per use-case, but
practitioners face stacks containing dozens of compensating candidates.
Without an explicit procedure for application, the construct risks
remaining theoretical: invoked at the level of architecture review, but
absent from the operating discipline that distinguishes one layer from
the next. The procedure that follows operationalizes the two conditions
of §3, producing a verdict in a finite number of steps for any (layer,
use-case) pair under examination. It can be applied without adopting the
instantiation of §4.

\begin{verbatim}
Input:  layer L deployed for use-case U.
Output: verdict in {symptom, structural};
        defect D identified when symptom.

Step 1 — Compensation
  Identify the defect D of the underlying persistence model
  for which L compensates in use-case U.
  If no such D is identifiable, return structural.

Step 2 — Dissolution
  Consider a counterfactual model M' in which D does not arise.
  If L retains a role in M' for U, return structural.
  Otherwise return symptom with D.
\end{verbatim}

Three notes on application. First, the unit of analysis is the pair
(layer, use-case), not the layer alone. The same product may produce
different verdicts for different uses; the procedure must be run once
per use-case. Second, identifying the defect in Step 1 requires explicit
attribution to a property of the persistence model. \emph{``We use this
layer because it is faster''} does not name a defect; \emph{``we use
this layer because the model serializes writes against a shared store''}
does. The procedure depends on the discipline of this attribution.
Third, the discriminator table of §3 and the Redis case of §4.2 are
worked examples of the procedure applied to canonical and
product-specific layers respectively. Each row in those tables is the
output of the procedure for one (layer, use-case) pair.

The procedure renders verdicts; it does not prescribe action. A symptom
verdict does not require that the layer be removed in the current
system. It indicates that the layer's continued presence is structurally
tied to retaining the persistence model that produces the defect.
Practitioners may rationally retain a symptom --- for legacy
preservation, contractual constraint, team familiarity, migration cost
--- but they do so knowing that the layer carries a debt that another
model would not impose. The procedure's product is visibility, not
directive.

\subsection{6. The ladder of compensated
layers}\label{the-ladder-of-compensated-layers}

The seven canonical categories enumerated in §3 are not a list. Read in
order, they form a ladder of increasing scale: in-memory cache
(per-process), object-relational mapper (within-application), message
queue and distributed lock (between processes), application server
framework (the application stack), orchestration cluster (across
machines), and full operating system (the hosting substrate itself). The
order is not arbitrary; it follows the expanding radius of the defect
the layer compensates for. Each rung compensates for a deeper limitation
of the model; each rung dissolves under the alternative for the same
reason.

Three additional categories that §3 forward-referenced extend the
ladder.

Backup and disaster-recovery tooling --- point-in-time restore,
write-ahead log archiving, snapshot orchestration --- exists because
state in DBMS-centric architectures lives in mutable rows whose history
is not the primary artifact. Under the model, the journal is the backup,
and rehydration is the recovery operation; both reduce to a single
substrate property treated more fully in Paper 5.

The service mesh --- inter-service routing, discovery, traffic shaping,
mTLS, sidecar instrumentation --- exists because stateless application
servers, by design, do not retain coordination state internally. Under
the model, much of this coordination is internal: Reactions and Tell
handle deferred work and cross-actor dispatch as journal entries. The
parts that cross the operator's boundary (TLS to external partners,
traffic shaping for external clients) remain structural and continue to
belong in a mesh or comparable substrate.

The distributed-tracing and APM stack --- span propagation, trace
ingestion, request-flow visualization, latency aggregation --- requires
discrimination per use-case, treated in Table 4 below.

\textbf{Table 4.} Use-case-level discrimination for the
distributed-tracing and APM stack.

{\def\LTcaptype{none} % do not increment counter
\begin{longtable}[]{@{}
  >{\raggedright\arraybackslash}p{(\linewidth - 4\tabcolsep) * \real{0.3333}}
  >{\raggedright\arraybackslash}p{(\linewidth - 4\tabcolsep) * \real{0.3333}}
  >{\raggedright\arraybackslash}p{(\linewidth - 4\tabcolsep) * \real{0.3333}}@{}}
\toprule\noalign{}
\begin{minipage}[b]{\linewidth}\raggedright
Observability use-case
\end{minipage} & \begin{minipage}[b]{\linewidth}\raggedright
Category
\end{minipage} & \begin{minipage}[b]{\linewidth}\raggedright
Reason
\end{minipage} \\
\midrule\noalign{}
\endhead
\bottomrule\noalign{}
\endlastfoot
Distributed tracing to reconstruct what happened inside the system &
Symptom & The journal already records the trace; replay is
deterministic \\
Distributed tracing across boundaries with external systems & Structural
& The external system does not dissolve under any persistence-model
change \\
APM as \emph{``find the bug in this incident''} & Symptom
(predominantly) & Replay localizes the bug without requiring external
instrumentation \\
APM as continuous capacity, latency, and error-rate telemetry in
production & Structural & Resource consumption is a real-world property,
not a model artifact \\
\end{longtable}
}

The first and third rows describe symptoms; the second and fourth
describe structural requirements. Debug as a one-shot operation
dissolves more cleanly than continuous telemetry because the journal
substitutes directly for the debugger. Continuous telemetry is mixed:
the internal-trace component dissolves, but the
capacity-and-cross-system component remains. §8 treats this as a worked
example of the construct's discrimination at the boundary of its
applicability.

The construct's verdicts do not eliminate the products from a
deployment. The framework does not dissolve APM and OTel; it dissolves
the internal-tracing use-case while leaving APM and OTel plugged in at
the points where they instantiate a structural requirement. The same
observation applies to mesh products at external boundaries, and to
backup tooling for legacy systems retained alongside a journaled core.

The ladder is not closed. New layers will emerge to compensate for new
defects of DBMS-centric architectures --- feature-flag servers, schema
registries, secret stores, sidecars for auth and identity,
configuration-as-a-service products. Each new layer can be audited under
the procedure of §5 and assigned a verdict under the construct of §3.
The ladder is open upward, sideways, and into the future because the
underlying persistence model continues to generate compensations; the
procedure walks each new rung as it appears.

\subsection{7. At the limit: actor-native systems on embedded
hardware}\label{at-the-limit-actor-native-systems-on-embedded-hardware}

The procedure of §5 walks each rung of the ladder as it appears; §6
showed the ladder extends. The natural question is what remains when the
procedure of §5 removes every rung it marks as symptom. §7 names that
limit.

What remains is the domain artifact, the journal that records its
execution, and the minimal substrate that runs them. The substrate
retains what is irreducibly needed: a runtime, journal storage, network
I/O, and a scheduling loop. Not required are the cluster orchestrator,
the application server framework, the object-relational mapper, the
external cache, the distributed coordinator, and the general-purpose
operating system in its hosting form. The domain artifact is the
application; the substrate is what is required to run it; everything
else has been audited and removed under the construct.

The same artifact runs in four deployment shapes without architectural
modification. In a multi-datacenter cluster with cross-region
replication, the journal and Tell handle topology; the artifact is
unchanged. In a Kubernetes red-black rolling deployment, the journal
mediates the release handoff. On a single server, the artifact runs
against local storage; nothing in the model requires the cluster form.
On minimal embedded hardware --- a point-of-sale terminal, a field
device, a branch agent --- the same artifact and journal run on a
substrate sized to the device. The difference between these shapes is
operational variation over the same architecture. The model contains no
components that mandate any particular scale.

The preceding is a claim about what the model permits, not a claim about
what has been operationally proven. The deployments of the framework's
home site run in conventional infrastructure across eight domain
installations; no point-of-sale terminal, no field device, and no
embedded agent is currently among them. The construct's prediction for
embedded deployment follows from the model's properties as documented in
§4 and §6; the demonstration that the prediction holds under field
conditions has not yet been performed. §7 makes a claim of architectural
permission, not a claim of operational proof. The distinction is offered
as a constraint on credibility, not as a hedge.

At the limit, only the artifact, the narrative of its execution, and the
minimal substrate remain --- nothing that smaller hardware cannot host.

\subsection{8. When the construct does not
apply}\label{when-the-construct-does-not-apply}

Sections 3 through 7 showed the construct in action: canonical layers in
§3, the Redis case in §4, the procedure in §5, the extended ladder in
§6, and the limit case in §7. Each of those sections assumed conditions
--- that the persistence model is identifiable, that use-cases are
decomposable, that a counterfactual model is conceivable. §8 treats the
cases in which one or more of those conditions friction against the
system being audited. The construct still applies; the verdicts become
coarser and the translation to action grows more complicated.

The canonical boundary case is observability. Table 4 (§6) shows that
the APM stack contains four use-cases with mixed verdicts: two symptom
(internal distributed tracing, one-shot debug) and two structural
(cross-system tracing, continuous capacity telemetry). Each verdict is
correctly assigned per use-case. But a single OTel deployment in a
production system serves all four use-cases simultaneously without
surgical separation. Removing the instrumentation that serves the
internal-tracing use-case affects the instrumentation that serves
cross-system tracing; the two share span propagation, collector
pipelines, and storage backends. The construct discriminates the
use-case, not the deployment. Operational reality does not admit the
same discrimination, and the practitioner retains the deployment intact
and carries the symptom use-cases with it.

Four limits constrain what the construct can be asked to do.

\textbf{Quantification.} The procedure returns a binary verdict per
use-case; it does not predict cost saved, performance gained, or
operational complexity reduced. Quantification requires instrumentation
distinct from the construct.

\textbf{Stack prescription.} A stack populated by symptoms can be
identified as suboptimal, but the construct does not prescribe an
alternative. The instantiation of §4 demonstrates that one exists; it
does not specify the unique alternative for any given system.

\textbf{Adoption prediction.} A symptom verdict does not inform the
difficulty of migration: organizational costs, contractual constraints,
vendor lock-in, and engineering skills carry independent weight.

\textbf{Opaque models.} The construct requires identifying the defect of
the underlying persistence model in Step 1 of §5. Systems whose
persistence model is opaque --- third-party SaaS exposing only an API,
proprietary platforms without published internals --- fall outside its
applicability.

Section 5's closing observation extends here: the visibility that the
construct renders is independent of the action the practitioner takes.
Symptoms may remain in deployments for legitimate reasons --- legacy
preservation, contractual obligation, organizational constraint,
migration cost. The construct documents the debt; it does not resolve
it.

The construct works best where the persistence model is explicit and
characterizable, where use-cases are decomposable, and where a
counterfactual model is conceivable. Where one of these conditions
fails, the construct still applies but its discriminations grow coarser.
Where it applies cleanly, what the construct surfaces is not a smaller
stack but a domain that the model permits to be modeled, libraried, and
evolved without being interleaved with compensating concerns. The
utility of the construct is proportional to the decomposability of the
system under examination.

Persistence is one axis along which representational pressure deforms a
domain --- the axis this paper has traced. There is at least one more.
When interaction is modeled before the domain --- when aggregates take
the shape of the screen rather than the screen rendering the actor's
output --- the same compensation pattern reappears in a different
register: UI components, view-models, request DTOs, and step-state
objects accumulate around a domain that no longer fits its own
substrate.

The construct of accidental category applies there with the same force
it applied to infrastructural layers and to the server-role. The actor's
existing output boundary --- the primitives by which it emits to its
invoker, agnostic of who consumes --- already factors transport out of
the domain. This axis is not developed here.

\subsection{9. Extension: from layers to
roles}\label{extension-from-layers-to-roles}

The construct of §3 discriminates among \emph{layers} of infrastructure
--- components that occupy positions in a stack and compensate for the
persistence model from those positions. Architectural \emph{roles} ---
the entity that accepts authoritative writes, the master in replication,
the bootstrap issuer, the orchestrator --- admit the same discrimination
under the same two conditions. A role is an accidental category when its
necessity is traceable to a specific representational choice (what nodes
replicate to share state) and the role loses justification when that
choice is replaced.

The same construct therefore applies beyond layers. Under a substrate in
which nodes replicate programs rather than data, state, or operations,
the server-role does not survive as a structural requirement. Its
dissolution is not an architectural objective; it is a structural
consequence of the same representational choice that dissolves the
layers described in §6.

The two analyses operate at different levels of abstraction --- first on
layers, then on the roles that those layers exist to serve --- without
modifying the construct itself. The discrimination scales upward
unchanged.

\subsection{Code provenance}\label{code-provenance}

Source-code references in this paper resolve against the public
Puppeteer repository at commit
\href{https://github.com/alvaroNCubo/puppeteer/tree/2f31f9674a5de816bdf1bf9d8360ff218a02e4da}{\texttt{2f31f96}}
(2026-05-18). The snapshot is archived in Software Heritage under the
following persistent identifier:

\begin{verbatim}
swh:1:dir:10e7e6bad7eb77b6c2e406762026177f95c3ae92;
  origin=https://github.com/alvaroNCubo/puppeteer;
  anchor=swh:1:rev:2f31f9674a5de816bdf1bf9d8360ff218a02e4da
\end{verbatim}

Inline references of the form \texttt{file.cs:NN} (e.g.,
\texttt{ActorHandler.cs:38}) resolve against this snapshot. A reader can
construct a per-file SWHID by adding the qualifiers
\texttt{;path=\textless{}path\textgreater{};lines=\textless{}NN\textgreater{}}
to the directory SWHID above. Future commits to the repository may
renumber lines; the SWHID preserves the cited state independently of any
future change to the repository or its hosting.

\subsection{Acknowledgments}\label{acknowledgments}

The author used large language models (including Claude and ChatGPT) as
editorial assistants for language refinement, structural feedback, and
literature navigation. All original ideas, terminology, theoretical
constructs, and technical content presented in this work are solely the
author's.

\begin{center}\rule{0.5\linewidth}{0.5pt}\end{center}

\subsection{References}\label{references}

Fowler, M. (2005). Event sourcing. Retrieved from
https://martinfowler.com/eaaDev/EventSourcing.html

Helland, P. (2007). Life beyond distributed transactions: An apostate's
opinion. \emph{Proceedings of the 3rd Biennial Conference on Innovative
Data Systems Research (CIDR)}, 132--141.

Helland, P. (2016). Immutability changes everything.
\emph{Communications of the ACM}, 59(1), 64--70.

Hevner, A. R., March, S. T., Park, J., \& Ram, S. (2004). Design science
in information systems research. \emph{MIS Quarterly}, 28(1), 75--105.

Hewitt, C., Bishop, P., \& Steiger, R. (1973). A universal modular ACTOR
formalism for artificial intelligence. \emph{Proceedings of the 3rd
International Joint Conference on Artificial Intelligence (IJCAI)},
235--245.

Kleppmann, M. (2017). \emph{Designing data-intensive applications: The
big ideas behind reliable, scalable, and maintainable systems}. O'Reilly
Media.

Rivera, A. (2026a). Anti-porous architecture: a unified design principle
for CQRS + Actor + Event-Sourcing systems. \emph{Puppeteer Papers
Series}, Paper 1.
https://github.com/alvaroNCubo/puppeteer-papers/blob/main/01-anti-porosity.md

Rivera, A. (2026b). Program--value separability: the structural
precondition for compilation, caching, and dense journaling in a DSL
runtime. \emph{Puppeteer Papers Series}, Paper 2.
https://github.com/alvaroNCubo/puppeteer-papers/blob/main/02-program-value-separability.md

Rivera, A. (2026c). Reactions and the partition: opt-in eventual
consistency in actor-native systems. \emph{Puppeteer Papers Series},
Paper 3.
https://github.com/alvaroNCubo/puppeteer-papers/blob/main/03-reactions-and-partition.md

Rivera, A. (2026d). Preserving semantic continuity across actors: a
tell-based approach without orchestration. \emph{Puppeteer Papers
Series}, Paper 4.
https://github.com/alvaroNCubo/puppeteer-papers/blob/main/04-cross-actor-continuity.md

Rivera, A. (2026e). The journal as substrate: unifying deployment,
replication, backup, and offline operation in distributed systems.
\emph{Puppeteer Papers Series}, Paper 5.
https://github.com/alvaroNCubo/puppeteer-papers/blob/main/05-substrate-operations.md

Rivera, A. (2026g). The server is not a structural requirement:
identifying accidental architectural roles under journaled programs.
\emph{Puppeteer Papers Series}, Paper 7.
https://github.com/alvaroNCubo/puppeteer-papers/blob/main/07-server-as-accidental-category.md

Stonebraker, M., \& Çetintemel, U. (2005). ``One size fits all'': An
idea whose time has come and gone. \emph{Proceedings of the 21st
International Conference on Data Engineering (ICDE)}, 2--11.

Vernon, V. (2015). \emph{Reactive messaging patterns with the actor
model: Applications and integration in Scala and Akka}. Addison-Wesley.

Young, G. (2010). CQRS documents. Retrieved from
https://cqrs.files.wordpress.com/2010/11/cqrs\_documents.pdf

\end{document}
